\documentclass[11pt]{article}
\usepackage{setspace}
\usepackage{geometry}
\geometry{margin=1in}
\setstretch{1.15}

\begin{document}

% --- Applicant Information ---
\noindent
\textbf{Zhongli Wang} \\
Phone: +49 160 9541 6015 \\
Email: nkwangzhongli@gmail.com \\
Website: https://sites.google.com/view/zhongli-wang \\
Bielefeld, Germany \\

\vspace{0.5cm}

% --- Recipient Information ---
\noindent
\textbf{Prof.\ Lukas Altermatt and Prof.\ Cyril Monnet} \\
Department of Economics \\
University of Bern \\

\vspace{0.4cm}

% --- Subject Line ---
\noindent
\textbf{Application -- SNSF-funded Postdoctoral Position: Financial Stability and Digital Currencies} \\

\vspace{0.3cm}

\noindent Dear Professor Altermatt and Professor Monnet,\\[0.2cm]

I am writing to apply for the SNSF-funded postdoctoral position in the Department of Economics at the University of Bern, associated with the projects \textit{``Financial Stability in an Era of Digital Currencies''} and \textit{``The Next-Generation Monetary System: Institutional Design of Interoperable and Programmable CBDCs.''} I am currently completing my Ph.D.\ in Economics at Bielefeld University and Universit\`a Cattolica del Sacro Cuore (Marie Sk\l{}odowska--Curie EPOC programme). My research specializes in computational and quantitative economic modeling---combining dynamic simulation, structural estimation, and data-driven methods---to study strategic behavior and policy-relevant counterfactuals in complex, digital environments. After reading your work, I am particularly excited by the projects' focus on oligopolistic competition in banking and on modeling digital currencies and monetary infrastructures, which I see as a close and natural extension of my own research interests.

My dissertation, \textit{``Strategic Behavior and Economic Dynamics in ICT Markets: An Agent-Based Framework,''} studies strategic interaction among firms and heterogeneous users in settings characterized by network effects, congestion, and boundedly rational decision rules. Methodologically, the dissertation develops computational models in which agents respond to incentives and information frictions over time; empirically, it relies on simulation-based econometrics to discipline model mechanisms and evaluate counterfactual policies. While the application domain is digital service markets, the underlying problems---strategic competition with frictions, dynamic adjustment, and heterogeneity---map closely onto modern monetary and financial systems, especially when digital currencies and payment innovations reshape participation, market structure, and stability.
Alongside my dissertation, I have coauthored projects that strengthen my fit through demonstrated strengths in data-intensive empirical work and research transparency. 
\textit{Polytope Fraud Theory} applies supervised machine learning to detect structured anomalies in firm-level financial statements, reflecting my ability to manage large datasets, engineer features, validate models, and communicate results with interpretability. A second project develops agentic workflows that integrate large-language-model agents into empirical research pipelines, emphasizing reproducibility, documentation, and scalable processing of complex data sources.

I am especially interested in extending my computational approach to the banking and monetary setting in two directions aligned with your projects. First, I would like to study how \emph{market structure and competition among banks or payment intermediaries} shape stability and welfare when new digital money instruments are introduced. Oligopolistic competition, product differentiation, and endogenous pricing of deposit-like claims create rich incentives for balance-sheet adjustment and liquidity management, and I am keen to contribute to modeling and calibrating these mechanisms in environments where digital money changes outside options, switching costs, and run incentives. 
Second, I am motivated to work on \emph{CBDC design and interoperability}, where institutional details---such as access rules, remuneration, programmability, and interoperability standards---can generate non-trivial equilibrium and stability effects. My prior work on network interactions and congestion provides a useful lens for understanding interoperability externalities, adoption dynamics, and potential ``capacity'' constraints in payments and settlement, as well as how these features interact with heterogeneity in users and intermediaries.

My technical background fits the projects' emphasis on theoretical modeling and programming. I have extensive experience implementing large-scale simulations and estimation pipelines in Python, including parallel computing for computationally intensive models and systematic sensitivity analysis for policy counterfactuals. 
% In my dissertation, I also employ surrogate modeling (e.g.\ neural-network approximations) to accelerate calibration and to explore high-dimensional parameter spaces. These tools are well suited for developing and calibrating quantitative models of digital monetary systems, where identifying key mechanisms and producing policy-relevant quantitative results often requires repeated solution and simulation under alternative design regimes. In addition, 
I have worked with large datasets and machine-learning methods in related projects, which positions me well to engage with empirical validation or data construction if the projects incorporate micro evidence or institutional data.

Beyond technical fit, I am drawn to the University of Bern because of its strong monetary economics environment and the opportunity to connect with the Swiss National Bank and the Study Center Gerzensee. I value research settings that combine rigorous modeling with policy relevance, and the SNSF project structure---including the expectation of independent work alongside close collaboration with the project leaders---matches my preferred way of working. I would be excited to contribute to model development, calibration, and paper writing, while also developing independent research that complements the central goals of the grants.

% Enclosed are my cover letter, CV, job market paper, and reference contacts, as requested. 
Thank you so much for considering my application! I would be delighted to discuss how my computational modeling background and my interest in banking competition and digital currency design can contribute to your research agenda at the University of Bern.

\vspace{6mm}

\noindent Sincerely,\\[12mm]
\noindent Zhongli Wang

\end{document}
